% Immutable cited-input supplement. Target proofs remain embargoed.
% Target: cdp:st:C.2.RAT.9; tag 01CL.
% Frozen source SHA256 74497bfe327fae25c9cdcb17334b3bf42a9793a9df7c0a4404d97568f2f760f0.
% CompendiumRef is resolved against the cited S6_09 source edition, not silently identified with the current series statement.

% Exact cited input: series:II / cdp:st:A.5.19
\begin{proposition}[Ordinary-triple images have arithmetic genus at most one]\hypertarget{tag-01B9}{}\label{cdp:st:A.5.19}


Let $f:Y\to C_0$ be a proper holomorphic map from a smooth
compact complex threefold with $b_2(Y)=0$, with connected
fibres and smooth general fibre a complex two-torus.
Let $mS$ be a full fibre with irreducible reduced support,
and let $D\subset\operatorname{Sing}S$ be an irreducible
curve with generic ordinary-triple plane transversal.
Then $D$ is a local complete intersection in $Y$ and
$p_a(D)\le1$. The possibilities for its underlying curve
are a smooth rational curve, a rational curve with one
ordinary node or cusp, and a smooth elliptic curve.
In the elliptic case the transversal type is constant.

When $D$ is singular rational, every component of the
normalized conductor curve dominating $D$ is rational,
and its map to $D^\nu\simeq\PP^1$ has degree one.
Consequently, in the irrational ruled setup of \ref{cdp:st:A.5.18}
with the generic transversal replaced by an ordinary
triple point, the image of the noncontracted horizontal
section is either smooth rational or smooth elliptic.
In the elliptic case the horizontal cover is unramified
of degree three, the normalization is smooth minimal
elliptic ruled, and the fibre has Euler number zero.
In the rational case the horizontal cover has degree
two or three; these cases are not excluded by this
statement.
\end{proposition}

% Exact cited input: series:II / cdp:st:A.5.16
\begin{proposition}[Rational normalization bounds with the exact constant-sheaf map]\hypertarget{tag-01B6}{}\label{cdp:st:A.5.16}


Let $f:Y\to B$ be a proper holomorphic map from a smooth compact
complex threefold with $b_2(Y)=0$, with connected fibres and smooth
general fibre a two-torus. Let $mS$ be a full fibre with irreducible
nonnormal reduced support. Write $\eta:T\to S$ for normalization,
$\pi:Z\to T$ for its minimal resolution, and $\Gamma\subset S$
for the conductor-image scheme. Suppose that $Z$ is rational.
Put $Q=\eta_*\OO_T/\OO_S$ and $e=h^0(S,\omega_S)$.
Then $e\in\{0,1\}$, and
\[
\begin{gathered}
 R^1\pi_*\OO_Z=0,\qquad
 H^1(T;\mathbb C)=H^1(T,\OO_T)=H^2(T,\OO_T)=0,\\
 H^0(S,Q)\xrightarrow{\ \simeq\ }H^1(S,\OO_S),\qquad
 h^0(S,Q)=1+e,\quad h^1(S,Q)=e,\quad
 \chi(\Gamma,\OO_\Gamma)=-1.
\end{gathered}
\tag{A.RT.1}\hypertarget{tag-0071}{}\label{cdp:eq:A.RT.1}
\]
In particular $T$ has rational singularities. The coefficient map
\[
 H^1(S;\mathbb C)\longrightarrow H^1(S,\OO_S)
\tag{A.RT.2}\hypertarget{tag-007C}{}\label{cdp:eq:A.RT.2}
\]
is injective, so $b_1(S)\le1+e\le2$.

There is an effective integral divisor $D$ on $Z$ with
$\OO_Z(D)\simeq K_Z^{-1}$, whose strict-transform part is the
normalization conductor. It satisfies
\[
 h^0(D,\OO_D)=h^1(D,\OO_D)=1,
 \quad D_{\rm red}\text{ connected},
 \quad p_a(D_{\rm red})\le1.
\tag{A.RT.3}\hypertarget{tag-007F}{}\label{cdp:eq:A.RT.3}
\]
The conductor image is connected,
$p_a(\Gamma_{\rm red})\le2$, and
$e(\Gamma_{\rm red})\ge0$. Set-theoretically
$D_{\rm red}=\pi^{-1}(\eta^{-1}(\Gamma_{\rm red}))$.
Exceptional curves over points outside the conductor have
coefficient zero in $D$; they are not to be counted as conductor
branches. With $E_p=(\pi^{-1}(p))_{\rm red}$ the exact Euler
formula is
\[
 e(S)=e(Z)-e(D_{\rm red})+e(\Gamma_{\rm red})
       -\sum_{p\in\operatorname{Sing}T\setminus\eta^{-1}(\Gamma)}
                  \bigl(e(E_p)-1\bigr).
\tag{A.RT.4}\hypertarget{tag-007G}{}\label{cdp:eq:A.RT.4}
\]
In the torus smoothing of the statement, \ref{cdp:st:A.5.17} excludes
every point in this sum. Hence $T$ is smooth outside the
normalization conductor, and the sum is zero.
\end{proposition}

% Exact cited input: series:II / cdp:st:A.5.18
\begin{proposition}[Every horizontal tacnode image is elliptic]\hypertarget{tag-01B8}{}\label{cdp:st:A.5.18}


Let $f:Y\to C_0$ be a proper holomorphic map from a smooth
compact complex threefold with $b_2(Y)=0$, with connected
fibres and smooth general fibre a complex two-torus.
Let $mS$ be a full fibre with irreducible reduced support,
$\eta:T\to S$ its normalization, and $\pi:Z\to T$ a
resolution with an irrational ruling $p:Z\to B$.
Suppose every exceptional curve of $\pi$ is vertical and
the conductor divisor on $Z$ is $2C+V$, where $C$ is a
section of $p$ and $V$ is vertical. Let $D=\eta\pi(C)$.
Suppose the generic plane-curve transversal to $D$ in $S$ is an
ordinary tacnode.

Then $D$ is a smooth elliptic curve, $B=C$ is elliptic,
and $C\to D$ is an unramified double cover. The transversal
type along $D$ is everywhere the same ordinary tacnode.
The normalization $T$ is a smooth minimal ruled surface
over this elliptic curve. Its conductor is twice a
section, and $e(S)=0$. In particular its minimal
resolution is the identity.
In particular a rational image of
the horizontal tacnodal conductor cannot occur. The proof
allows vertical conductor and rational singularities of
$T$; it does not assume that the normalization map is
immersive at the special points under examination.
\end{proposition}

% Exact cited input: series:II / cdp:st:A.5.17
\begin{proposition}[Isolated singularities are excluded and suspension germs are transverse products]\hypertarget{tag-01B7}{}\label{cdp:st:A.5.17}


Let $f:Y\to C_0$ have the hypotheses of \ref{cdp:st:A.5.15}.
No local equation $f=u h^m$, with $u$ a holomorphic
unit and $m\ge1$, can have a reduced equation $h$
with an isolated critical point at the origin.
Suppose its local equation at a point of a fibre is
\[
             f=(x^a-g(y,z))^m,
             \qquad a\ge2,\quad m\ge1,             \tag{A.J.38}\hypertarget{tag-0069}{}\label{cdp:eq:A.J.38}
\]
where $g\ne0$ and $g(0,0)=0$. Then the reduced
plane germ $\{g=0\}_{\rm red}$ is smooth. If
$b$ is its multiplicity in $\operatorname{div}(g)$,
there are local holomorphic coordinates in which
$f=(x^a-y^b)^m$.
Consequently, at every point where the full fibre
is a uniform multiple of its reduced support and
the reduced equation has a nonzero Hessian, its
critical surface germ is a constant
transverse $A_{b-1}$ product. A nonordinary change
of transverse type at a point with nonzero Hessian
is excluded.
\end{proposition}

% Exact cited input: series:II / cdp:st:A.5.15
\begin{proposition}[Hirzebruch germs cannot occur in a torus degeneration]\hypertarget{tag-01B5}{}\label{cdp:st:A.5.15}


Let $f:Y\to C_0$ be a proper holomorphic map from a
smooth complex threefold to a smooth curve, with smooth
nearby fibres complex two-tori. There is no point of
a fibre at which, in holomorphic ambient coordinates
$(x,y,z)$ and a coordinate on the base, the map has
the equation
\[
 f(x,y,z)=u(x,y,z)\bigl(x^a-y^bz^c\bigr)^m,
 \quad a\ge2,\quad b,c,m\ge1,                                    \tag{A.J.33}\hypertarget{tag-0062}{}\label{cdp:eq:A.J.33}
\]
where $u$ is a holomorphic unit. No Betti-number
assumption on $Y$ is needed. This includes the pinch
point, every germ $x^2-y^{2k}z^{2r+1}$ with
$k\ge1$, $r\ge0$, and the reducible germ
$x^2-y^2z^2$.
\end{proposition}

